\chapter{今後の計画}\chaplab{plan}
本章では，最終報告に向けた今後の活動計画について述べる．はじめに班分けによる活動体制を説明し，次にモデルを用いたシミュレーションによる検証プロセス，最後に運行案の提案までの流れを述べる．
\bunseki{山本貫太}

\section{活動体制}
今後のモデルを用いたシミュレーションおよび運行案の策定に向けて，モデル班とデータ班の2つの班に分かれて活動を行う．次項にて，それぞれの班が担う具体的な役割と活動内容について述べる．
\bunseki{山本貫太}

\subsection{モデル班}
モデル班は，現実のバス運行状況を再現するためのモデル構築，および快速化案の導入効果を定量的に測るための評価指標の策定を担当する．

バスの運行をシミュレーションするためには，乗降客数の変動や車両の遅延といった要因を数式として正確に表現する必要がある．そのため，本班ではまず基盤となる数理モデルの概念や背後にある理論について理解し，その作成方法に関する体系的な学習から着手する．学習を通じて得た知見をもとに，現状の運行を再現するためにはどのようなアプローチが最適であるか，モデルの手法を検討・選定して設計を進めていく．

手法の選定と並行して，現実のバス運行において発生する事象のうち「何を定式化し，何を削ぎ落とすか」の定義を行う．現実のバス運行は，交通状況や個々の乗客の行動など無数の要因が絡み合う複雑な現象である．シミュレーション上でこれを機能させるためには，すべての事象をそのまま再現するのではなく，混雑緩和という目的に対して「何が本質的な要因か」を見極める抽象化のプロセスが不可欠となる．影響の微小な不要な要素を捨象し，乗客の乗降やバスの走行といった中核となるプロセスのみを抽出して数式に落とし込む．モデルの内部構造や各パラメータがシミュレーション結果に与える影響を正しく把握した上で定式化を行い，独自のモデルを構築していく．

さらに，次段階の検証において快速化案の導入効果を客観的かつ定量的に評価するための指標づくりも，モデル班の役割となる．本プロジェクトでは，主要な課題である混雑という抽象的な状態をいかにして数値として定義するかが重要となる．積み残しの人数や最大乗車率，平均待ち時間，所要時間，遅延時間のほか，通過停留所の利用者に生じる待ち時間の増加など，どのような要素を基準として設定すべきかを検討し，評価指標の策定を行う．
\bunseki{山本貫太}

\subsection{データ班}
データ班は，モデル班が構築するシミュレーションを実際に稼働させ，現実の事象を再現するために必要なデータの収集や整理をする役割を担う．

数理モデルで現実のバス運行状況を再現するためには，実データが必要不可欠である．具体的には，公共交通の標準データフォーマットであるGTFS（General Transit Feed Specification）データを活用して，路線や停留所の位置情報，および基準となる運行ダイヤを構築する．さらに，各停留所における時間帯別の乗客数や，車両ごとの遅延状況を示す実際の運行データなどが必要であると考えられる．しかし，本プロジェクトで活用を想定しているこれらのGTFSデータや乗客数，実際の運行データは，一般に公開されているオープンデータではない．そのため，検証に必要な情報を得るために運行主体である函館バス株式会社との連携が必須となる．実態に即した精度の高い前提条件を設定するべく，同社への直接的な相談やデータ提供の依頼を実施する計画である．

本プロジェクトが目指すのは，適度に抽象化された単純なモデルを通じて「どのような運行方法が効果的か」という原理を見出すことである．したがって，データ整理においては，モデル班が抽出した「本質的な要素」に合致する粒度へとデータを丸めていく．また，データ提供の協議と併せて，実際の運行管理において現場が抱えている課題感や，特定の時間帯・区間における交通事情などを直接ヒアリングする．これにより，数値データだけでは読み取れない定性的な情報も補完し，シミュレーションの精度をより高めていく．これらの活動を通じて，モデル班が設計するモデルに合致したデータを供給し，次段階のシミュレーション検証フェーズへの移行を図る．
\bunseki{山本貫太}

\section{モデルを用いたシミュレーション}
本節では，前節で述べたモデル班とデータ班の準備成果を統合し，実際にシミュレーションを稼働させて仮説を検証していくプロセスについて述べる．

実際のバス路線において運行方法を変更し，複数の条件を繰り返し試すことは，日常的にバスを利用する学生や一般利用者，バス事業者に影響を与えるため容易ではない．そこで本プロジェクトでは，バスの運行や利用者の乗降などを数理モデルとして表現し，仮想的に運行条件を変更することで，それぞれの条件が混雑や遅延に与える影響を同じ基準のもとで比較する．ここでの数理モデルの役割は，現実のバス運行を完全に再現することではなく，解決策の効果を評価するために必要な要素を抽出し，異なる運行条件を比較できるようにすることである．そのため，モデルから得られる結果については，設定した仮定や使用したデータを踏まえて解釈し，現実の状況と比較しながら妥当性を検討する．

モデル化の対象は，函館駅前から未来大学までを結ぶ一本の路線とする．この路線は停留所がほぼ一列に並ぶため一次元的に扱いやすく，また，快速化と2周運行という提案がこの区間内で完結して表現できるためである．想定しているシミュレーションでは，時間帯ごとの利用者数，停留所ごとの乗降人数，バスの発車時刻，停留所間の所要時間，バスの定員などを入力として用いる．各停留所に乗客が到着し，バスが時刻表に従って移動し，停車する停留所で乗客が乗降する流れを再現し，バスの定員を超えた場合には，乗れなかった乗客が次の便を待つものとして扱う．これにより，現実に近い形で混雑や積み残しを表現できると考えている．

検証は大きく分けて現状の再現と快速化の検証の二段階で構成し，その過程において反復的な改善ループを実施する．

本プロジェクトが構築するのは事象の本質を抽出した単純なモデルであるため，シミュレーション結果と現実の観測データとの間にいくらかの乖離が生じるのは当然の前提とする．したがって，現実の運行を秒単位や人単位で緻密に模倣しようと追求するのではなく，混雑の発生メカニズムを正しく捉えられているかどうかの確認に主眼を置く．細かな数値の一致を求めて過度な修正を繰り返すと，かえってモデルから本質が失われ，汎用性のないものになってしまう危険性がある．そのため，現実とシミュレーションの間に「根本的な原理のズレ」が生じた場合にのみ，モデルの前提条件やデータの丸め方に立ち戻って修正を加える．このように，過度な適合に陥るのを防ぎつつ，検証の土台として必要十分な妥当性を確保していく．

精度の高い現状の再現ができた後，第二段階として快速化の検証へと移行する．現状を再現できているモデルに対して，特定の停留所を通過させる快速バスを導入する．また，快速化によって短縮できた時間を利用して同じ車両が混雑区間をもう一度運行するといった新たな条件を導入し，シミュレーションを再実行する．その際，快速便を運行する時間帯や本数，停車する停留所の組み合わせを変えた複数のパターンを検証する．その結果生じる運行状況の変化について，モデル班が事前に定義した評価指標を用い，導入前後での比較・分析を行う．これにより，快速化の導入が混雑緩和や所要時間の短縮にどの程度の効果をもたらすのかを定量的に明らかにし，次節における具体的な運行案の提案へ向けた客観的な根拠を提示する．
\bunseki{山本貫太}

\section{運行案の提案}
本節では，シミュレーションによる検証結果をもとに，具体的な解決策を提示するプロセスについて述べる．

初めに，検証結果の総合的な評価を行う．前節で行なった現状の運行を再現したシミュレーションおよび複数の快速化案に関するシミュレーション結果を比較・分析する．モデル班が策定した評価指標に基づき，どの快速化のパターンが効果的に混雑を削減し，所要時間の短縮に寄与できるかを客観的かつ定量的に評価する．

続いて，その評価結果をもとに具体的な運行案の策定を行う．最も高い効果が実証された快速化案をベースに，「学生の利便性」，「バス事業者の運行効率」，「一般利用者への影響」の3つを考慮し，実際の路線に適用可能な運行案として集約する．

最終的に，策定された運行案は，検証の根拠となる客観的なデータとともに，実際の運行主体である函館バス株式会社へ提案する計画である．本プロジェクトの最終目標は，数理モデルを用いた検証を通じて混雑しないバスの走らせ方を見つけ出し，未来大学の学生の毎日の移動をもっと快適にすることにある．実効性のある運行案を提案することで，「満員で乗れないかもしれない」，「授業に遅れるかもしれない」といった学生の日常的な不安を払拭し，時間に余裕のある通学環境の実現に貢献することを目指す．
\bunseki{山本貫太}
